\scanpage{122}
\margin{\textcircled{3}}
Nous en savons assez (finitude, Künneth, dualité
globale) pour faire le Lefschetz ``éculé et \unsure{poisseux|poussiéreux}'',
en redescendant sur terre. Il semble que pour les
courbes, on récupère par des arguments d'\unsure{immersion}
(sans résolution des singularités) des Lefschetz plus
\unsure{astucieux}, pour faisceaux singuliers $\rightsquigarrow$
rationalité des fonctions $L$.

Nous pouvons aussi, au lieu de cela,
continuer ces grandes \unsure{fondations|fonctions}, en l'occurrence :
\uuline{th. de bidualité}, \inserted{(= \uline{dualité locale}),}
\uline{Lefschetz}-\uline{Cartier}, \deleted{et}
\deleted{\unsure{encore les fonctions $L$}}
\deleted{\unsure{où utiliser le} \illegible{1}}
\deleted{\unsure{alors}} ce nouveau pavé.

Étant parti dans un formulaire de dualité
général, je préfère en compléter les lignes essentielles
avant de \deleted{\unsure{enfin}} me mettre aux applications de
détail.
\mspage{121}
